\documentclass[11pt]{article}

\usepackage[T1]{fontenc}
\usepackage[utf8]{inputenc}
\usepackage{amsmath}
\usepackage{amssymb}
\usepackage{array}
\usepackage{booktabs}
\usepackage{enumitem}
\usepackage[margin=1in]{geometry}
\usepackage{hyperref}
\usepackage{listings}
\usepackage{xcolor}

\definecolor{codebackground}{RGB}{246,247,249}
\definecolor{codeframe}{RGB}{210,214,220}
\definecolor{linkblue}{RGB}{25,82,135}

\hypersetup{
  colorlinks=true,
  linkcolor=linkblue,
  citecolor=linkblue,
  urlcolor=linkblue,
  pdftitle={Pinndorama Solver Architecture and Numerical Implementation},
  pdfauthor={Pinndorama contributors}
}

\lstdefinestyle{pinndorama}{
  basicstyle=\ttfamily\small,
  backgroundcolor=\color{codebackground},
  frame=single,
  rulecolor=\color{codeframe},
  columns=fullflexible,
  keepspaces=true,
  showstringspaces=false,
  breaklines=true,
  upquote=true
}
\lstset{style=pinndorama}

\newcommand{\code}[1]{\texttt{\detokenize{#1}}}
\newcommand{\mean}{\operatorname{mean}}
\newcommand{\JRH}{J R_H}
\newcommand{\dd}{\mathrm{d}}

\title{Pinndorama Solver Architecture and Numerical Implementation}
\author{Pinndorama contributors}
\date{July 2026}

\begin{document}

\maketitle

\begin{abstract}
Pinndorama is a collection of physics-informed neural-network solvers for
nonlinear elliptic equations.  It contains a one-dimensional manufactured
spherical problem and three conformally flat binary-puncture applications:
fixed axisymmetric, equal-spin parametric axisymmetric, and fully
three-dimensional.  This document describes the active implementation as a
technical reference.  It connects the common package structure to the
mathematical residuals, boundary losses, collocation rules, optimizers,
checkpoint format, evaluation tools, and vendored dependencies.  It documents
what the code computes; it does not introduce accuracy, convergence, or
performance claims, and no trained checkpoints or numerical results are
assumed to be present.
\end{abstract}

\tableofcontents

\section{Purpose and scope}

The installable Python sources live under
\path{src/pinndorama}.  Four solver packages are maintained independently:

\begin{itemize}[nosep]
  \item \path{solvers/toy_problem_1d}: a non-parametric spherical manufactured
        problem with configurable outer boundary treatment;
  \item \path{solvers/punctures_2d}: a fixed-physics axisymmetric puncture
        solver;
  \item \path{solvers/punctures_2d_parametric}: an axisymmetric puncture solver
        whose third input is a physical equal-spin parameter; and
  \item \path{solvers/punctures_3d}: the full three-dimensional puncture
        solver.
\end{itemize}

The packages deliberately do not inherit from a shared solver class.  They
have parallel organization and share numerical conventions, but each retains
its own configuration validation, model, loss, optimizer adapter, checkpoint
code, and evaluator.  This keeps the scientific inputs to each application
explicit.  Shared repository resources are resolved centrally by
\path{src/pinndorama/_paths.py}; the shared SSBroyden loader is
\path{src/pinndorama/_ssbroyden.py}.

Throughout the code, \code{xx0}, \code{xx1}, and \code{xx2} denote native
curvilinear coordinates.  Their collocation counts are \code{Nxx0},
\code{Nxx1}, and \code{Nxx2}.  These are PINN sample counts, not
finite-difference dimensions, and they contain no ghost zones.  Physical
coordinates are denoted by $(x,y,z)$ or by spherical $(r,\vartheta,\varphi)$.

\section{Package organization and common data flow}

\subsection{Directory structure}

The solver subtree has the following conceptual structure.  Files shown in
the common column occur in every solver; the final group records
application-specific additions.

\begin{lstlisting}
src/pinndorama/
|-- _paths.py                 repository resource locations
|-- _ssbroyden.py             shared authenticated optimizer loader
|-- reproducibility/          dispatch, reference metadata, comparisons
`-- solvers/
    |-- toy_problem_1d/
    |-- punctures_2d/
    |-- punctures_2d_parametric/
    `-- punctures_3d/

Each solver package:
    config.py                  strict TOML schema and effective settings
    coordinates.py             maps and collocation/boundary samples
    model.py                   float64 MLP and output transform
    loss.py                    automatic derivatives and total objective
    train.py                   Adam, SSBroyden, logging, checkpoints
    ssbroyden_trainer.py       full-loss adapter to the shared runtime
    checkpoint.py              atomic versioned NPZ persistence
    evaluate.py                explicit checkpoint evaluation

Puncture packages additionally:
    nrpy_expression_builder.py symbolic reference metric and source setup
    generated_expressions.py   cached NumPy/JAX numerical callables

The 1D package additionally:
    equation.py                manufactured equation and exact solution
    plot_error.py              explicit evaluation-table error plot
\end{lstlisting}

The root \path{configs/} directory contains explicit TOML inputs grouped by
solver name.  The Python package does not silently select a default research
configuration.  The standalone C applications under
\path{reference_solvers/nrpyelliptic/} are not installed as Python modules.
Pinned third-party sources live under \path{vendor/}.

\subsection{End-to-end data flow}

A training run follows the same broad sequence in all four packages:

\begin{enumerate}
  \item Parse and strictly validate a user-supplied TOML file.  Apply only the
        supported command-line overrides for collocation counts, seed, and
        optimizer-stage lengths.
  \item Construct cell-centered interior collocation points and separate,
        structured samples for each boundary or regularity term.
  \item Construct coordinate quantities.  For puncture problems, build cached
        NRPy/SymPy expressions and precompute the fixed background and
        Bowen--York source arrays on the interior samples.
  \item Initialize or resume a float64 multilayer perceptron (MLP).
  \item Use JAX automatic differentiation to obtain the first and diagonal
        second derivatives of the network output in native coordinates.
  \item Evaluate the mean-square interior residual and the independent
        boundary losses, then form their configured weighted sum.
  \item Apply an Optax Adam warm-up followed by the shared, vendored
        SSBroyden quasi-Newton refinement.
  \item Write periodic and final atomic NPZ checkpoints with JSON metadata.
        Evaluation later loads an explicit checkpoint and an explicit
        coordinate file; no run directory is searched implicitly.
\end{enumerate}

Fixed-shape interior and boundary arrays are constructed before JIT
compilation.  Network evaluation, differentiation, loss assembly, and Adam
steps use JAX arrays.  Scalar residual functions are vectorized with
\code{jax.vmap}; repeated call graphs are compiled with \code{jax.jit} or the
Equinox filtered JIT used by the Optimistix wrapper.

\subsection{Configuration-first interface}

The canonical training interface is

\begin{lstlisting}[language=bash]
python -m pinndorama.solvers.<solver>.train \
  --config CONFIG.toml --output-dir OUTPUT_DIR \
  [--resume-checkpoint CHECKPOINT.npz]
\end{lstlisting}

The command line may override native resolution, the random seed, Adam
length, and SSBroyden length.  Physics, coordinate geometry, architecture,
activation, output ansatz, loss operators, boundary sampling policy,
optimizer tolerances, and learning rates remain in TOML.  Configuration keys
are validated exactly; old or unknown names are not translated silently.

\section{Common neural-network representation}

\subsection{MLP definition}

Let $q\in\mathbb{R}^{n_{\rm in}}$ be the model input.  With $L$ dense layers,
the code evaluates
\begin{align}
  h^{(0)} &= q,\\
  h^{(\ell)} &= \tanh\!\left(W^{(\ell)}h^{(\ell-1)}+b^{(\ell)}\right),
       &&1\leq \ell<L,\\
  f_{\rm raw}(q) &= W^{(L)}h^{(L-1)}+b^{(L)}.
\end{align}
The output layer is linear and scalar.  Dense weights are stored in Equinox
layout, $(n_{\rm out},n_{\rm in})$, even though the forward pass is written
as a row-major batch product with $W^{T}$.  Checked configurations use He
initialization,
\begin{equation}
  W_{ij}^{(\ell)}\sim
  \mathcal{N}\!\left(0,\frac{2}{n_{\rm in}^{(\ell)}}\right),
  \qquad b_i^{(\ell)}=0,
\end{equation}
and enable JAX 64-bit mode before creating parameters or evaluating losses.

The usual \code{w40d4} architecture has four hidden layers of width 40:
\begin{equation}
  [n_{\rm in},40,40,40,40,1].
\end{equation}
The fixed 2D family also contains a checked \code{w60d10} architecture with
ten hidden layers of width 60.

\subsection{Smooth inverse-radius ansatz}

All solvers turn the raw output into the trained field using
\begin{equation}
  u(q)=\frac{f_{\rm raw}(q)}{\sqrt{1+r(q)^2}}.
  \label{eq:output-transform}
\end{equation}
This transform supplies smooth behavior at finite $r$ and an inverse-radius
falloff at large radius without fixing the leading asymptotic coefficient.
Boundary losses still constrain derivatives and asymptotic structure; the
ansatz does not make those losses identically zero.

The model inputs are
\begin{center}
\begin{tabular}{lll}
\toprule
Solver & Input vector & Input dimension \\
\midrule
1D toy & $(xx0)$ & 1 \\
Fixed 2D & $(xx0,xx1)$ & 2 \\
Parametric 2D & $(xx0,xx1,\texttt{equal\_spin\_sz})$ & 3 \\
Full 3D & $(xx0,xx1,xx2)$ & 3 \\
\bottomrule
\end{tabular}
\end{center}

The mass $m$ in the toy problem is configuration-selected physics and is not
a network input.  Conversely, \code{equal_spin_sz} is deliberately an input
to the parametric network.

\section{Native coordinate systems}

\subsection{SinhSpherical map for the 1D solver}

The toy solver directly implements the NRPy SinhSpherical radial map
\begin{equation}
 r(xx0)=\mathrm{AMPL}\,
 \frac{\sinh(xx0/\mathrm{SINHW})}{\sinh(1/\mathrm{SINHW})},
 \qquad 0\leq xx0\leq1.
 \label{eq:sinh-spherical}
\end{equation}
The analytic map derivatives are
\begin{align}
 r_{,0} &=
 \frac{\mathrm{AMPL}}{\mathrm{SINHW}\sinh(1/\mathrm{SINHW})}
 \cosh(xx0/\mathrm{SINHW}),\\
 r_{,00} &=%
 \frac{\mathrm{AMPL}}{\mathrm{SINHW}^{2}\sinh(1/\mathrm{SINHW})}
 \sinh(xx0/\mathrm{SINHW}).
\end{align}
JAX differentiates $u$ with respect to $xx0$; the code then computes
\begin{equation}
 u_r=\frac{u_{,0}}{r_{,0}},\qquad
 u_{rr}=\frac{u_{,00}}{r_{,0}^{2}}
         -\frac{u_{,0}r_{,00}}{r_{,0}^{3}}.
 \label{eq:radial-chain-rule}
\end{equation}

\subsection{SinhSymTP map for the puncture solvers}

For the puncture applications, define
\begin{equation}
 A(xx0)=\mathrm{AMAX}\,
 \frac{\sinh(xx0/\mathrm{SINHWAA})}
      {\sinh(1/\mathrm{SINHWAA})}.
\end{equation}
With $b=\texttt{bScale}$, the native-to-Cartesian map is
\begin{align}
 x &= A\sin(xx1)\cos(xx2),\\
 y &= A\sin(xx1)\sin(xx2),\\
 z &= \sqrt{A^2+b^2}\cos(xx1).
 \label{eq:sinh-symtp-cart}
\end{align}
The spherical radius used in the output transform and Robin term is
\begin{equation}
 r=\sqrt{A^2+b^2\cos^2(xx1)}.
\end{equation}
SinhSymTP is orthogonal.  The implemented reference-metric scale factors are
\begin{align}
 h_0 &= A_{,0}
       \frac{\sqrt{A^2+b^2\sin^2(xx1)}}{\sqrt{A^2+b^2}},\\
 h_1 &= \sqrt{A^2+b^2\sin^2(xx1)},\\
 h_2 &= A\sin(xx1),
\end{align}
so
\begin{equation}
 J=\sqrt{\det\hat\gamma}=h_0h_1h_2,
 \qquad \hat\gamma^{ii}=h_i^{-2}.
 \label{eq:sinh-symtp-jacobian}
\end{equation}

The outer condition needs $u_r$ at fixed spherical polar angle, not simply a
native $xx0$ derivative.  The coordinate modules analytically form the
Jacobian of $(r,\vartheta)$ with respect to $(xx0,xx1)$ and invert the
two-dimensional gradient transformation:
\begin{equation}
 u_r=
 \frac{u_{,0}\vartheta_{,1}-\vartheta_{,0}u_{,1}}
      {r_{,0}\vartheta_{,1}-\vartheta_{,0}r_{,1}}.
 \label{eq:spherical-radial-derivative}
\end{equation}
The $xx2$ derivative does not enter this transformation because $r$ and
$\vartheta$ are independent of azimuth.

\section{Puncture Hamiltonian residual}

\subsection{Equation and source fields}

The three puncture solvers represent a correction $u$ to the singular
conformal-factor background $\psi_{\rm bg}$.  In the notation exposed by the
NRPy source builder, define
\begin{equation}
 \mathcal{A}=\texttt{ADD\_times\_AUU}.
\end{equation}
The conformally flat Hamiltonian residual is
\begin{equation}
 R_H=\hat\nabla^2u+
 \frac{\mathcal{A}}{8}(\psi_{\rm bg}+u)^{-7}.
 \label{eq:hamiltonian-residual}
\end{equation}
The bare masses, momenta, spins, and z-axis separation \code{zPunc} select the
Bowen--York two-puncture source.  No transverse puncture-separation input is
accepted.

\subsection{Jacobian-regularized divergence form}

Training does not use the raw expression in
Eq.~\eqref{eq:hamiltonian-residual}.  For an orthogonal reference metric it
uses the algebraically equivalent divergence form multiplied by $J$:
\begin{equation}
 \JRH=
 \sum_{i\in\mathcal{D}}
 \left[
   \partial_i\!\left(J\hat\gamma^{ii}\right)u_{,i}
   +J\hat\gamma^{ii}u_{,ii}
 \right]
 +J\frac{\mathcal{A}}{8}(\psi_{\rm bg}+u)^{-7}.
 \label{eq:jrh}
\end{equation}
Here $\mathcal{D}=\{0,1\}$ for the axisymmetric solvers and
$\mathcal{D}=\{0,1,2\}$ in 3D.  Multiplication by $J$ removes coordinate
singular denominators algebraically.  The tests verify the identity
$\JRH=J R_H$ away from coordinate singularities.

The symbolic builders use NRPy to construct $J$, $\psi_{\rm bg}$, and
$\mathcal{A}$ with a stable ordered symbol interface.  The 2D builders set
NRPy's symmetry axis to direction 2 before constructing expressions, so all
$xx2$ derivatives vanish symbolically rather than being discarded after
generation.  SymPy \code{lambdify} produces cached NumPy and JAX callables.
For a fixed collocation array, $\psi_{\rm bg}$ and $\mathcal{A}$ are evaluated
once before optimization.  The derivatives of $u$ are still recomputed from
the current network parameters by JAX at every loss evaluation.

\section{Training and optimization}

\subsection{Adam warm-up}

The first stage uses Optax Adam with a cosine-decay learning-rate schedule
from the configured initial rate to the configured minimum fraction.  Each
gradient leaf is clipped elementwise to $[-c,c]$, where $c$ is the configured
clip value.  The clipping is not a global gradient-norm rescaling.  The entire
weighted loss, including every active boundary term, is differentiated in one
\code{jax.value_and_grad} call.

The checked schedules are summarized below.

\begin{center}
\begin{tabular}{lrrrr}
\toprule
Configuration & Adam steps & Initial rate & SSBroyden steps & $r_{\rm tol}=a_{\rm tol}$ \\
\midrule
T001 & 1,500 & $10^{-4}$ & 10,000 & $10^{-12}$ \\
C002/C004 & 1,500 & $10^{-4}$ & 50,000 & $10^{-12}$ \\
P001 & 1,500 & $10^{-4}$ & 50,000 & $10^{-12}$ \\
C04 & 10,000 & $10^{-3}$ & 50,000 & $10^{-12}$ \\
C05 & 0 & --- & 50,000 & $10^{-12}$ \\
\bottomrule
\end{tabular}
\end{center}

\subsection{Shared SSBroyden refinement}

After Adam, each solver passes the same complete scalar loss to the shared
SSBroyden runtime.  SSBroyden is a self-scaled member of the Broyden
quasi-Newton family~\cite{bioli2026}.  It stores an inverse-Hessian
approximation as a PyTree linear operator, uses a Newton descent with a
Cholesky linear solver, and uses Optimistix's Zoom line search by default.
The self-scaling parameters $\theta_k$ and $\tau_k$ are recomputed on accepted
updates.  In the notation of the vendored implementation, its inverse-Hessian
update has the form
\begin{equation}
 H_{k+1}=
 \frac{1}{\tau_k}\left[
 H_k-\frac{H_ky_ky_k^TH_k}{y_k^TH_ky_k}
 +\phi_k(y_k^TH_ky_k)v_kv_k^T
 \right]+\rho_ks_ks_k^T,
 \label{eq:ssbroyden}
\end{equation}
where
\begin{equation}
 \rho_k=(y_k^Ts_k)^{-1},\qquad
 v_k=\rho_ks_k-\frac{H_ky_k}{y_k^TH_ky_k},\qquad
 \phi_k=\frac{1-\theta_k}{1+a_k\theta_k}.
\end{equation}
The relative and absolute tolerances are passed to Optimistix's Cauchy-style
termination test; they are not direct target values for the PINN loss.

The wrapper JIT-compiles the step and termination functions, counts accepted
line-search steps, and supplies the accepted state's stored loss to callbacks.
Each solver separately tracks the lowest accepted loss and returns those
parameters, rather than assuming the last trial state is best.  Periodic
checkpoints therefore store the current best accepted parameters.

\subsection{Continuation semantics}

Continuation restores only model parameters and scientific metadata.  Adam
moments, SSBroyden inverse-Hessian state, and line-search state are not stored.
A resumed stage creates a new optimizer state around the loaded model.  This
is why collocation and optimizer settings may change between parent and child
runs.  The child metadata records the SHA-256 digest of the parent checkpoint.

\section{The one-dimensional spherical toy problem}

\subsection{Equation and manufactured solution}

The mass $m>0$ is selected by the TOML file.  The manufactured solution and
singular background are
\begin{equation}
 u_{\rm exact}(r)=\frac{2m(1+r)}{r^2+2r+2},
 \qquad
 \psi(r)=1+\frac{m}{2r}.
 \label{eq:toy-exact}
\end{equation}
The equation is
\begin{equation}
 R=u_{rr}+\frac{2}{r}u_r+
 \frac{1}{8}A(r)[\psi(r)+u(r)]^{-7}=0,
 \label{eq:toy-pde}
\end{equation}
with
\begin{equation}
 A(r)=
 \frac{m(r^2+6r+6)
 [2r(r^2+2r+2)+m(5r^2+6r+2)]^7}
 {4r^7(r^2+2r+2)^{10}}.
\end{equation}
Although $A$ and $\psi$ are separately singular at the origin, the nonlinear
combination is finite.  The loss evaluates the exactly equivalent expression
\begin{equation}
 N(r,u,m)=\frac{1}{8}A(\psi+u)^{-7}
 =\frac{4m(r^2+6r+6)
 [2r(r^2+2r+2)+m(5r^2+6r+2)]^7}
 {(r^2+2r+2)^{10}[m+2r(1+u)]^7}.
 \label{eq:toy-stable-source}
\end{equation}
The denominator is evaluated as written; the implementation does not clip or
regularize it numerically.

\subsection{Interior and origin terms}

The interior points are the $Nxx0$ cell centers of $[xx0_{\min},xx0_{\max}]$.
The trained residual is regularized by one power of physical radius:
\begin{equation}
 rR=r u_{rr}+2u_r+rN(r,u,m).
 \label{eq:toy-regularized-residual}
\end{equation}
The coordinate origin is not an interior cell center.  It is evaluated
separately at $xx0=0$, where the regularity residual is
\begin{equation}
 B_{\rm origin}=u_r(0).
\end{equation}

\subsection{Outer boundary operators}

The TOML independently selects one operator and one enforcement region.  The
\emph{first-order Robin} residual is
\begin{equation}
 B_1(r)=(1+r)(ru_r+u).
 \label{eq:first-order-robin}
\end{equation}
It annihilates a leading $C/r$ profile.  The factor $(1+r)$ is part of the
implemented residual and keeps the constraint scaled at large radius.

The \emph{second-order Robin} residual is
\begin{equation}
 B_2(r)=r^3\left(u_r+\frac{u}{r}\right)-2m.
 \label{eq:second-order-robin}
\end{equation}
It incorporates the known subleading term in
\begin{equation}
 u(r)=\frac{C}{r}-\frac{2m}{r^2}+\mathcal{O}(r^{-3}),
\end{equation}
without supplying $C$ to the network.

The two enforcement regions are:
\begin{description}[style=nextline]
  \item[\code{endpoint}]
    Evaluate the selected residual only at the explicit endpoint $xx0=1$.
    The configuration must omit \code{outer_boundary_r_min}.
  \item[\code{radial_band}]
    Select every interior cell center whose mapped radius strictly satisfies
    $r>r_{\min}^{\rm outer}$, then append the explicit $xx0=1$ endpoint.  Here
    $r_{\min}^{\rm outer}$ is the value of
    \code{outer_boundary_r_min}; it must satisfy
    $0<r_{\min}^{\rm outer}<\mathrm{AMPL}$.
\end{description}
Thus the supported strategies are $B_1$ at the endpoint, $B_1$ on a radial
band, $B_2$ at the endpoint, and $B_2$ on a radial band.  Exactly one strategy
is active in a run.  The checked T001 configuration uses $B_2$ on the strict
$r>10^3$ band plus the endpoint.

\subsection{Total toy loss}

Let $\mathcal{I}$ be the cell-centered interior set and $\mathcal{B}$ the
selected outer set.  With $B$ equal to $B_1$ or $B_2$, the implemented loss is
\begin{equation}
 \mathcal{L}_{\rm toy}=
 w_R\mean_{q\in\mathcal{I}}[rR(q)]^2
 +w_0[u_r(0)]^2
 +w_B\mean_{q\in\mathcal{B}}[B(q)]^2.
 \label{eq:toy-loss}
\end{equation}
T001 sets all three weights to one, uses 256 cell-centered points, and uses a
$[1,40,40,40,40,1]$ network.  The exact solution in
Eq.~\eqref{eq:toy-exact} is not part of
Eq.~\eqref{eq:toy-loss}; it is used only by tests and evaluation output.

\section{Fixed axisymmetric punctures}

\subsection{Inputs and collocation}

The fixed 2D network consumes $(xx0,xx1)$.  Every physical source component is
fixed by the selected TOML.  Interior samples form a flattened tensor product
of $Nxx0$ and $Nxx1$ cell centers.  The checked fixed-2D configurations are:

\begin{center}
\begin{tabular}{llll}
\toprule
Configuration & Collocation & Architecture & Role \\
\midrule
C002 \code{w40d4} & $384\times192$ & $[2,40,40,40,40,1]$ & baseline \\
C002 $70\times70$ & $70\times70$ & $[2,40,40,40,40,1]$ & lighter run \\
C002 \code{w60d10} & $384\times192$ & 10 layers of width 60 & wider/deeper \\
C004 square & $256\times256$ & $[2,40,40,40,40,1]$ & square grid \\
\bottomrule
\end{tabular}
\end{center}

\subsection{Boundary sample sets}

Theta scalar parity is imposed with structured pairs at both polar
boundaries.  For samples $(xx0,\delta)$,
\begin{align}
 E_{\rm lower} &=u(xx0,-\delta)-u(xx0,\delta),\\
 E_{\rm upper} &=u(xx0,\pi+\delta)-u(xx0,\pi-\delta),\\
 \mathcal{L}_{\theta} &=\frac12\left[
 \mean E_{\rm lower}^2+\mean E_{\rm upper}^2\right].
 \label{eq:2d-parity}
\end{align}
The $xx0$ and $\delta$ arrays are cell-centered and use $Nxx0\times Nxx1$
samples; $\delta$ spans the configured near-axis interval.

The outer samples set $xx0=xx0_{\max}$ and use $Nxx1$ cell centers in the
polar coordinate.  Their scaled Robin residual is
\begin{equation}
 B_{\infty}=(1+r)(ru_r+u),
 \qquad
 \mathcal{L}_{\rm out}=\mean B_{\infty}^2,
 \label{eq:puncture-robin}
\end{equation}
where $u_r$ is computed by Eq.~\eqref{eq:spherical-radial-derivative}.

\subsection{Total fixed-2D loss}

With $\mathcal{L}_{J}=\mean(\JRH)^2$, the fixed-2D loss is
\begin{equation}
 \mathcal{L}_{\rm 2D}=
 w_J\mathcal{L}_{J}+w_{\theta}\mathcal{L}_{\theta}
 +w_{\rm out}\mathcal{L}_{\rm out}.
 \label{eq:fixed-2d-loss}
\end{equation}
All checked configurations set these weights to one.  Adam runs for 1,500
steps at an initial rate of $10^{-4}$ and SSBroyden is allowed up to 50,000
accepted steps.

\section{Equal-spin parametric axisymmetric punctures}

The parametric solver retains the same 2D SinhSymTP mapping, axisymmetric
residual in Eq.~\eqref{eq:jrh}, scalar-parity pairs in
Eq.~\eqref{eq:2d-parity}, and outer residual in
Eq.~\eqref{eq:puncture-robin}.  Its distinction is the third model input
\begin{equation}
 s=\texttt{equal\_spin\_sz}.
\end{equation}
This is the raw physical Bowen--York $S_z$ component, not a dimensionless spin
$\chi$.  For every source evaluation,
\begin{equation}
 S_{0z}=S_{1z}=s.
\end{equation}
The remaining source components are fixed by the TOML.

P001 uses 80 cell centers in $s\in[-0.2,0.2]$ and a
$128\times64$ native spatial grid.  The tensor product therefore contains
$128\times64\times80$ interior residual points.  Each spatial parity and
outer-boundary sample is likewise repeated at all 80 spin cell centers.  The
parameter axis is not called \code{Nxx2}, because it is not a spatial
SinhSymTP coordinate.

The source builder substitutes $S_{0z}=S_{1z}=s$ before evaluating
$\psi_{\rm bg}$ and $\mathcal{A}$.  JAX differentiates the model only with
respect to $xx0$ and $xx1$ for the PDE; no derivative with respect to $s$
appears in the elliptic residual.  With the checked
\code{regularized_residual_scale} equal to one, the total loss is
\begin{equation}
 \mathcal{L}_{\rm param}=
 w_J\mean\left(\frac{\JRH}{s_R}\right)^2
 +w_{\theta}\mathcal{L}_{\theta}
 +w_{\rm out}\mathcal{L}_{\rm out},
 \qquad s_R=1.
 \label{eq:parametric-loss}
\end{equation}
All checked weights are one.  P001 uses the
$[3,40,40,40,40,1]$ network, 1,500 Adam steps at $10^{-4}$, and up to 50,000
SSBroyden steps.

At evaluation time the user supplies one fixed $s$ in the configured range.
The returned field is a member of the trained parametric family, whereas a
fixed-2D checkpoint represents only the single physical source encoded by its
configuration.

\section{Full three-dimensional punctures}

\subsection{Interior residual and collocation}

The 3D MLP consumes $(xx0,xx1,xx2)$, and Eq.~\eqref{eq:jrh} contains all three
diagonal divergence terms.  The interior grid is the flattened tensor product
of cell-centered arrays in each native direction.

C04 uses $256\times256\times16$ samples with uniform $xx0$ cell centers over
$[0,1]$.  It performs 10,000 Adam steps at initial rate $10^{-3}$ followed by
up to 50,000 SSBroyden steps.  C05 must resume a compatible C04 checkpoint.
It keeps $Nxx0=256$ but changes the radial-like sampling to two zones:
\begin{itemize}[nosep]
  \item 236 cell centers on $[0,0.5]$;
  \item 20 cell centers on $[0.5,1]$.
\end{itemize}
C05 disables Adam and applies up to 50,000 new SSBroyden steps to the C04
parameters on the changed collocation set.

\subsection{Axis regularity}

Axis samples form a structured product of $xx0$, a small positive $\delta$,
and azimuth $\varphi$.  The code defines the wrapped antipodal angle
\begin{equation}
 \varphi_* = \operatorname{wrap}_{[-\pi,\pi)}(\varphi+\pi)
 =[(\varphi+2\pi)\bmod 2\pi]-\pi.
\end{equation}
The lower and upper residuals are
\begin{align}
 E_{0} &=u(xx0,-\delta,\varphi_*)-u(xx0,\delta,\varphi),\\
 E_{\pi} &=u(xx0,\pi+\delta,\varphi_*)
            -u(xx0,\pi-\delta,\varphi),\\
 \mathcal{L}_{\rm axis}&=\frac12
 \left[\mean E_0^2+\mean E_{\pi}^2\right].
 \label{eq:3d-axis-loss}
\end{align}
The checked configurations sample $xx0$ from
\code{boundary_xx0_min} to the outer endpoint, $\delta$ from
\code{xx1_axis_delta_min} to \code{xx1_axis_delta_max}, and azimuth without a
duplicated endpoint.

\subsection{Azimuthal periodicity and outer boundary}

For a structured $(xx0,xx1)$ set away from the polar coordinate singularity,
\begin{equation}
 \mathcal{L}_{\varphi}=\mean
 \left[u(xx0,xx1,xx2_{\max})-u(xx0,xx1,xx2_{\min})\right]^2.
 \label{eq:phi-periodicity}
\end{equation}
The outer samples set $xx0=xx0_{\max}$ and form a product of polar points
away from the axes and azimuthal points without a duplicated endpoint.  They
use the same $B_{\infty}$ and $\mathcal{L}_{\rm out}$ as
Eq.~\eqref{eq:puncture-robin}.

\subsection{Total 3D loss}

The implemented objective is
\begin{equation}
 \mathcal{L}_{\rm 3D}=\mean(\JRH)^2
 +w_{\rm axis}\mathcal{L}_{\rm axis}
 +w_{\varphi}\mathcal{L}_{\varphi}
 +w_{\rm out}\mathcal{L}_{\rm out}.
 \label{eq:3d-loss}
\end{equation}
C04 and C05 set the three boundary weights to one.  The interior coefficient
is fixed to one in the 3D loss implementation.

\section{Checkpoints and evaluation}

\subsection{Atomic NPZ format}

Every solver writes \path{checkpoint_final.npz}; stage checkpoints use names
such as \path{checkpoint_adam_00001000.npz} and
\path{checkpoint_ssbroyden_00001000.npz}.  Checkpoint schema 1 contains
\begin{itemize}[nosep]
  \item float64 arrays \code{leaf_000}, \code{leaf_001}, and so on, in stable
        $W_0,b_0,W_1,b_1,\ldots$ order; and
  \item \code{metadata_json}, a one-dimensional \code{uint8} array containing
        UTF-8 JSON.
\end{itemize}
Loading uses \code{allow_pickle=False}.  Array count, shape, dtype, finiteness,
and contiguous leaf names are validated.  Saving writes a temporary file in
the destination directory, flushes it, and atomically replaces the requested
path.  Pickle checkpoint compatibility is intentionally absent.

Resume validation treats dtype, geometry, physics, architecture, and ansatz as
immutable.  The parametric solver additionally treats its declared parameter
range as immutable.  Loss choices, collocation, and optimizer-stage settings
are metadata but may change on continuation.  The child records
\code{parent_checkpoint_sha256}, allowing an external workflow to reconstruct
the parameter lineage.

\subsection{Explicit evaluation}

The shared dispatcher is
\begin{lstlisting}[language=bash]
python -m pinndorama.reproducibility.evaluate SOLVER \
  --config CONFIG.toml --checkpoint CHECKPOINT.npz \
  --coords COORDS --output OUTPUT
\end{lstlisting}
where \code{SOLVER} is one of
\begin{center}
\code{toy-problem-1d},\quad
\code{punctures-2d},\quad
\code{punctures-2d-parametric},\quad
\code{punctures-3d}.
\end{center}
The parametric form also requires \code{--equal-spin-sz}.  The dispatcher
invokes the selected solver evaluator in a subprocess using the same Python
interpreter.

Coordinate files contain native coordinates, either as text columns or as an
NPZ with exactly the expected named arrays.  The toy evaluator writes
$xx0,r,u_{\rm NN},u_{\rm exact}$ and pointwise errors.  Puncture evaluators
write the native point, mapped Cartesian coordinates, and $u_{\rm NN}$.
Every checkpoint, coordinate file, output, and physics configuration is
explicit; no cluster path or bundled result is assumed.

\section{Vendored dependencies}

\subsection{NRPy symbolic construction}

\path{vendor/nrpy/} is an ordinary-file snapshot of the upstream NRPy
repository~\cite{ruchlin2018}.  Its provenance records
\begin{itemize}[nosep]
  \item upstream URL \url{https://github.com/nrpy/nrpy};
  \item commit
        \code{4d16f0e35ce90a50b87508088d2dfcf744aa65b9};
  \item upstream release \code{2.2026.6}; and
  \item snapshot date 2026-07-10.
\end{itemize}
The upstream Git metadata is excluded, but the source snapshot is otherwise
kept separate from Pinndorama adaptations.

The central path resolver exposes \code{NRPY_ROOT}.  Each puncture symbolic
builder prepends that directory before importing NRPy reference-metric and
NRPyElliptic source modules.  The fixed-2D builder additionally checks
\code{nrpy.__file__} and rejects a module outside the vendored root.  The
parametric and 3D builders rely on the prepended central path, while the
symbolic-builder tests run each package in a fresh process and verify the
pinned source-symbol contract.  This behavior is intentionally stated
precisely: NRPy loading is not authenticated by the same common loader used
for SSBroyden.

The toy problem does not import NRPy at runtime.  Its small SinhSpherical map
is implemented directly and tested against the pinned formula.  The puncture
solvers use NRPy only to construct symbolic expressions and numerical
callables; training does not download code or regenerate C applications.

\subsection{Shared SSBroyden and Optimistix runtime}

\path{vendor/ssbroyden/} contains a single runtime shared by all four neural
solvers:
\begin{itemize}[nosep]
  \item \path{ssbroyden_family.py}, from
        \url{https://github.com/IvanBioli/ssbroyden_optimistix} commit
        \code{4c87785c68f0fec6b09000f474daef76fb181eea};
  \item \path{optimistix_wrapper.py}, the accepted-step manual loop; and
  \item \path{optimistix/}, from the required Optimistix fork, declaring
        version \code{0.1.0}.  Its pinned commit is
        \path{8cd4931713658f8dfe4423ead6f11b348b675540}.
\end{itemize}
These are ordinary files rather than submodules, so the optimizer remains
available if the upstream repositories disappear.

The loader prepends both vendor roots, invalidates import caches, imports
Optimistix, the wrapper, and the SSBroyden family, and then verifies every
module's \code{__file__} lies beneath the expected vendor directory.  It also
requires \code{optimistix.__version__ == "0.1.0"}, rejecting an ambient
installation or version mismatch.  Optimistix reports that constant directly
instead of querying installed distribution metadata, so no separate editable
Optimistix installation is required.

Local adaptations recorded in \path{VENDORED_FROM.md} are the normalized
\path{ssbroyden_family.py} filename, accepted-state loss reuse in the wrapper,
the deterministic Optimistix version constant, and formatting-only changes.
\path{SHA256SUMS} covers the active runtime files and is checked by tests.  The
nested Optimistix snapshot retains its Apache-2.0 license.  The SSBroyden
add-on had no separate license at the pinned commit; written redistribution
permission is pending, and Pinndorama does not assign one.

\section{NRPyElliptic reference workflow}

Standalone NRPyElliptic C applications are stored under
\path{reference_solvers/nrpyelliptic/}.  Its axisymmetric subtree contains the
axisymmetric reference application.  Its three-dimensional subtree contains
the standard 3D application, the canonical comparison reader, and a
parametric-network warm-start variant.  These programs build in their own
directories and are not part of the Python installation.  NRPyElliptic's
hyperbolic-relaxation method is described in Ref.~\cite{assumpcao2022}.

The canonical reader has a path-neutral interface
\begin{lstlisting}[language=bash]
nrpyell_reader --binary FILE --coords FILE --output FILE
\end{lstlisting}
and consumes Cartesian $x,y,z$ points.  It reads the versioned
\code{NRPYELL3} binary format and interpolates the $u$ field with a
tensor-product stencil using nine source points in each native direction.  It
is the comparison reader for both axisymmetric and 3D neural solutions.

An NRPy solution header does not encode every Bowen--York source parameter.
The Python workflow therefore requires a reference TOML stored beside the
binary.  That file records the binary name and SHA-256, all SinhSymTP geometry,
and all source parameters.  The comparison rejects a changed digest, geometry
inconsistent with the binary header, or physics inconsistent with the solver
TOML.

Given neural values $u_{\rm NN}^{(a)}$, interpolated reference values
$u_{\rm NRPy}^{(a)}$, and explicit nonnegative volume weights $w_a$, the
reported metric is
\begin{equation}
 \epsilon_{L_2}=
 \sqrt{
 \frac{\sum_a w_a\left(u_{\rm NN}^{(a)}-u_{\rm NRPy}^{(a)}\right)^2}
      {\sum_a w_a\left(u_{\rm NRPy}^{(a)}\right)^2}}
 .
 \label{eq:relative-l2}
\end{equation}
The arrays must have identical nonzero size, all entries must be finite, at
least one weight must be positive, and the weighted reference norm must be
nonzero.

The optional warm-start C application accepts \code{initial_guess = zero} or
\code{initial_guess = parametric_nn}; restart data takes precedence.  No
weights are embedded.  The \code{make nn-header CHECKPOINT=...} target exports
a user-supplied P001-compatible $[3,40,40,40,40,1]$ checkpoint.  Before the
OpenMP initialization loop, the C code validates geometry, fixed source
parameters, equal raw spins, and the exported spin range.

\section{Representative commands}

The following commands illustrate the supported interfaces.  They require an
editable Pinndorama installation and write only to the explicit destinations.

\subsection*{Training}

\begin{lstlisting}[language=bash]
python -m pinndorama.solvers.toy_problem_1d.train \
  --config configs/toy_problem_1d/T001_w40d4.toml \
  --output-dir runs/T001_w40d4

python -m pinndorama.solvers.punctures_2d.train \
  --config configs/punctures_2d/C002_w40d4.toml \
  --output-dir runs/C002_w40d4

python -m pinndorama.solvers.punctures_2d_parametric.train \
  --config configs/punctures_2d_parametric/P001_w40d4.toml \
  --output-dir runs/P001_w40d4

python -m pinndorama.solvers.punctures_3d.train \
  --config configs/punctures_3d/C04.toml \
  --output-dir runs/C04

python -m pinndorama.solvers.punctures_3d.train \
  --config configs/punctures_3d/C05.toml \
  --output-dir runs/C05 \
  --resume-checkpoint runs/C04/checkpoint_final.npz
\end{lstlisting}

\subsection*{Evaluation and comparison}

\begin{lstlisting}[language=bash]
python -m pinndorama.reproducibility.evaluate punctures-3d \
  --config configs/punctures_3d/C04.toml \
  --checkpoint runs/C04/checkpoint_final.npz \
  --coords /path/to/native_xx.txt \
  --output /path/to/checkpoint_values.txt

make -C reference_solvers/nrpyelliptic/\
READER_nrpyelliptic_conformally_flat

python -m pinndorama.reproducibility.create_reference \
  --binary /path/to/NRPYELL_solution.bin \
  --solver-config configs/punctures_3d/C04.toml \
  --output /path/to/NRPYELL_solution.reference.toml

python -m pinndorama.reproducibility.compare \
  --nrpy-binary /path/to/NRPYELL_solution.bin \
  --reference-config /path/to/NRPYELL_solution.reference.toml \
  --solver-config configs/punctures_3d/C04.toml \
  --coords /path/to/cartesian_xyz.txt \
  --nn-values /path/to/checkpoint_values.txt \
  --volume-weights /path/to/quadrature_weights.txt \
  --output /path/to/comparison.json
\end{lstlisting}

\section{Implementation boundaries}

This reference describes the active code paths and checked configurations.
The four solvers share conventions but remain separate implementations; a
change in one package does not automatically alter the others.  Cluster
launchers, machine-specific paths, implicit artifact searches, and bundled
paper-scale results are outside these modules.

\begin{thebibliography}{9}

\bibitem{ruchlin2018}
I. Ruchlin, Z. B. Etienne, and T. W. Baumgarte,
``SENR/NRPy+: Numerical relativity in singular curvilinear coordinate
systems,''
\emph{Physical Review D} \textbf{97}, 064036 (2018).
\href{https://doi.org/10.1103/PhysRevD.97.064036}
{doi:10.1103/PhysRevD.97.064036}.

\bibitem{assumpcao2022}
T. Assumpcao, L. R. Werneck, T. P. Jacques, and Z. B. Etienne,
``Fast hyperbolic relaxation elliptic solver for numerical relativity:
Conformally flat, binary puncture initial data,''
\emph{Physical Review D} \textbf{105}, 104037 (2022).
\href{https://doi.org/10.1103/PhysRevD.105.104037}
{doi:10.1103/PhysRevD.105.104037}.

\bibitem{bioli2026}
I. Bioli and M. Mendibe Abarrategi,
``Self-Scaled Broyden Family of Quasi-Newton Methods in JAX,''
arXiv:2603.10599 (2026).
\url{https://arxiv.org/abs/2603.10599}.

\bibitem{optimistix}
P. Kidger and contributors,
``Optimistix: nonlinear optimization in JAX,'' software project.
The Pinndorama runtime uses the fork and commit recorded in
\path{vendor/ssbroyden/VENDORED_FROM.md}.
\url{https://github.com/patrick-kidger/optimistix}.

\end{thebibliography}

\end{document}
